\section{Experimental Design}
\label{sec:experimental-design}

\subsection{Frozen experimental protocol}
\label{subsec:frozen-protocol}

The definitive experiment follows protocol \texttt{paper1-q3-v1}, frozen before
execution of the experimental matrix. It specifies the research questions,
hypotheses, configuration families, seeds, repetitions, resource limits,
statistical procedures, and treatment of incomplete executions.

The scientific baseline corresponds to release \texttt{v1.1.0-rc.1}. The frozen
toolchain consists of TLA+ Tools 1.7.4, Alloy 6.2.0 with SAT4J, Java 17,
Python 3.12, and GNU \texttt{/usr/bin/time -v} for time and memory
instrumentation.

The definitive design covers bounded property verification, scientific mutants,
bounded implementation-model trace conformance, and TLC fault profiles.
Configuration families were enumerated before definitive execution. No
configuration absent from the frozen matrix was added after observing outcomes.

\subsection{Configuration matrix}
\label{subsec:configuration-matrix}

The experimental matrix contains fourteen configuration families:

\begin{itemize}
\item six valid-model scalability families, with small, medium, and large
profiles for TLC and Alloy
\item two mutation families, one per formal tool
\item two trace-conformance families, valid and negative
\item four TLC fault-profile families covering normal execution, replay
pressure, timeout behavior, and insufficient quorum
\end{itemize}

These families expand through properties, mutants, repetitions, scenarios, and
seeds into 1272 scheduled tasks. The count includes warm-up and measured
executions where defined, while RQ3 contributes one replay execution per
case-seed pair.

Each task records protocol and configuration identifiers, tool, model kind,
profile, repetition or seed, resource limits, and provenance metadata.

\subsection{Bound profiles}
\label{subsec:bound-profiles}

Three ordered profiles define the principal bounded scalability configurations
in Table~\ref{tab:bound-profiles}.

\begin{table}[htbp]
\centering
\footnotesize
\setlength{\tabcolsep}{2pt}
\caption{Bound profiles used in the definitive experiment.}
\label{tab:bound-profiles}
\begin{tabular}{lrrrrrrr}
\hline
Profile &
Shards &
Transfers &
Validators &
\shortstack{TLC\\quorum} &
\shortstack{TLC receipt\\copies} &
\shortstack{Alloy Receipt\\scope} &
\shortstack{Alloy State\\scope} \\
\hline
Small  & 2 & 1 & 2 & 2 & 1 & 2 & 6  \\
Medium & 3 & 2 & 3 & 2 & 2 & 4 & 8  \\
Large  & 4 & 3 & 4 & 3 & 2 & 6 & 10 \\
\hline
\end{tabular}
\end{table}

Shard, transfer, and validator counts are shared across the profile records.
TLC quorum and receipt-copy columns are TLC parameters. Alloy instead uses
Receipt scopes 2, 4, and 6, Message scopes 10, 20, and 30, and State scopes
6, 8, and 10 for small, medium, and large. Its commit predicate requires at
least two votes in every profile, so TLC quorum values are not Alloy quorum
thresholds.

Profiles are experimental bundles, not one-factor perturbations: shared
configuration sizes and tool-specific bounds vary together. Cost changes are
therefore associated with the complete profile rather than attributed causally
to one parameter. The \texttt{catalog} profile is reserved for RQ3 scenarios.

\subsection{Mutation configurations}
\label{subsec:mutation-configurations}

Mutation experiments use the medium profile. Five scientific defect classes are
instantiated for each formal representation:

\begin{itemize}
\item removal of receipt replay protection
\item destination credit before valid receipt consumption
\item commit after an abort-like terminal decision
\item timeout without release of source funds
\item quorum bypass
\end{itemize}

Each mutant has a predefined target property. After two warm-ups, ten measured
repetitions are executed per mutant, and logical detection must remain
consistent across measured repetitions.

\subsection{Trace-conformance configurations}
\label{subsec:trace-configurations}

RQ3 uses a separate multiseed design with ten valid scenarios and ten negative
cases. Each case is evaluated under thirty deterministic seeds, producing
300 valid and 300 negative trace executions. The definitive seed sequence is
\texttt{2026001} through \texttt{2026030}.

RQ3 therefore contains 600 TLC replay executions. Each case-seed pair executes
once without warm-up because the primary outcome is deterministic
classification, not repeated performance estimation.

\subsection{Repetitions and execution order}
\label{subsec:repetitions}

Model-checking configurations used for time and memory measurement have two
warm-up repetitions, ten measured repetitions, serial execution, and exclusion
of warm-ups from primary statistical analysis. Warm-up observations remain in
the execution record but do not contribute to the reported cost summaries.

Logical outcomes are expected to remain consistent across repetitions.
Repetition characterizes elapsed time and memory rather than creating
independent logical property outcomes. No run is repeated because it produces a
counterexample, timeout, or other unfavorable scientific result.

\subsection{Resource limits}
\label{subsec:resource-limits}

Each measured model-checking execution uses:

\begin{itemize}
\item timeout: 1800 seconds
\item maximum resident-memory limit: 12288 MiB
\item maximum concurrent experimental executions: one
\item TLC workers: one
\item Alloy solver: SAT4J
\end{itemize}

Time and maximum resident memory are collected with
\texttt{/usr/bin/time -v}. Timeout and out-of-memory outcomes remain
scientifically meaningful scalability observations in the experimental record
and are not removed from the matrix.

\subsection{Execution environment}
\label{subsec:execution-environment}

The definitive timing and memory campaign runs on a dedicated native Linux host.
Virtualized environments, WSL, and shared continuous-integration runners are
excluded from primary timing and memory measurements. Continuous integration is
used for functional, structural-integrity, and reproducibility checks, but its
timings are not mixed with definitive performance measurements.

Recorded environment metadata includes operating system and kernel, processor
architecture and count, available memory, Java and Python versions, formal-tool
versions, repository commit, and virtualization status.

\subsection{Measured outcomes}
\label{subsec:measured-outcomes}

For TLC, recorded measurements include:

\begin{itemize}
\item property result
\item generated states
\item distinct states
\item exploration depth
\item elapsed time
\item maximum resident memory
\item process exit status
\item counterexample information when applicable
\end{itemize}

For Alloy, recorded measurements include:

\begin{itemize}
\item assertion result
\item solver and scope
\item counterexample information
\item solving and total elapsed time when available
\item maximum resident memory
\item process exit status
\end{itemize}

TLC state counts and Alloy scopes are not equivalent state-space units.

Trace-conformance records include case identifier, seed, case type, acceptance
or rejection, diagnostic agreement, abstract and concrete steps, action,
transfer identifier, elapsed time, memory, source commit, and input hashes.

\subsection{Incomplete executions and censoring}
\label{subsec:censoring}

Executions may terminate as completed runs, counterexamples, timeouts,
out-of-memory outcomes, tool errors, or instrumentation failures. Timeout and
out-of-memory outcomes are retained as censored observations without replacing
their elapsed-time or memory values by configured limits or synthetic values.

Unexpected tool exits remain tool errors. Instrumentation failures may be
retried only under the frozen retry policy and never on the basis of the
scientific result. This treatment preserves unfavorable outcomes and prevents
selective removal of executions that would weaken a hypothesis.

\subsection{Statistical analysis}
\label{subsec:statistical-analysis}

Time and memory distributions are summarized by median and interquartile range,
with minimum, maximum, and coefficient of variation also retained. A 95\%
bootstrap confidence interval for the median uses 10,000 resamples. Wilson 95\%
confidence intervals are used for proportions, including mutation detection and
trace classification.

RQ4 reports profile-level medians and within-tool growth ratios. Spearman rank
association describes monotonic association between ordered profile and cost
when sufficient completed profiles are available.

TLC and Alloy are not compared through absolute execution-time claims because
their modelling semantics and state-space representations differ. Such
measurements are not treated as evidence of intrinsic tool superiority. The
three ordered profiles provide descriptive bounded-cost evidence but do not
establish an asymptotic complexity class or a general nonlinear growth law.
